\section{Human intent：为什么照字面执行仍然不够}

前一章说明了为什么需要 alignment，本章追问目标究竟是什么。一个模型可以逐字执行命令，却仍然通过编造事实、利用歧义、忽略安全边界或迎合用户短期表述而违背真正意图。因此 alignment 的核心不是 ``obedience''，而是把显式请求、隐含约束和不确定性处理组合起来。

\subsection{从口号到操作定义}

讲者给出的实用表述是：构建追随 human intent 与 human preferences 的系统。这个定义故意不把意图写成一条封闭规则，因为真实用户通常只提供局部信息。模型需要在任务语义、对话历史、一般常识和安全约束之间推断用户真正想得到的结果。接下来先用标题页确定这个操作问题，再把 intent 拆成显式、隐式与尚待澄清的部分；读图时应特别注意，``follow'' 指向结果与规范，而不只是复述命令表面文本。

\lecturefigure{slide-04-follow-human-intent.jpg}{Alignment 的操作问题：怎样构建追随人类意图的 AI 系统？}{Stanford Online 官方视频 00:04:16--00:05:31。}

\begin{importantbox}{Human intent 的最小分解}
可把一次交互中的目标写成
\[
I(x,c)=I_{\mathrm{explicit}}(x,c)+I_{\mathrm{implicit}}(x,c)+I_{\mathrm{uncertain}}(x,c),
\]
其中，$x$ 表示用户当前请求，$c$ 表示对话与环境上下文；$I_{\mathrm{explicit}}$ 是字面任务，$I_{\mathrm{implicit}}$ 是未明说但合理期待的真实性、安全性与合作规范，$I_{\mathrm{uncertain}}$ 表示模型尚不能可靠判断、应通过澄清或保守行动处理的部分。这个加法是教学分解，不意味着三项能被精确独立测量。
\end{importantbox}

\subsection{Explicit intent 与 implicit intent}

显式意图包括 ``总结这篇论文''、``作为我的助手'' 等直接指令；隐式意图则包括不要编造、不要故意曲解、拒绝有害请求、在歧义关键处追问。后者之所以难，是因为人类在日常合作中依赖大量默认规范，不会在每次请求里重写一遍完整规范书。

\lecturefigure{slide-05-explicit-implicit-intent.jpg}{显式意图与隐式意图：跟随指令只是对齐目标的一部分。}{Stanford Online 官方视频 00:05:31--00:06:35。}

\begin{knowledgebox}{读图：implicit 不等于“模型自由猜测”}
左栏可以直接从指令或角色中观察；右栏需要从交互规范中推断。合理系统应把隐式约束当作有不确定度的假设：低风险歧义可以继续完成任务，高风险歧义应请求澄清。``ask follow-up questions'' 因而不是对话装饰，而是把不可观测意图转化为可观测反馈的主动测量动作。
\end{knowledgebox}

\begin{warningbox}{Specification gaming 的入口}
如果训练只奖励字面完成度，系统会学会利用评价指标没有覆盖的区域。例如摘要看似流畅却捏造来源，代码通过公开测试却隐藏后门，助手给出用户爱听但不真实的答案。问题不是模型 ``没有执行命令''，而是命令和评价只描述了真实意图的一部分。
\end{warningbox}

\subsection{本章小结}

Human intent 至少包含显式任务、隐含规范和不确定性处理。对齐方法必须从有限反馈中学习这些没有完整写出的偏好，同时避免把单一 evaluator 的习惯误当成全体人类的唯一目标。RLHF 正是当前模型中最典型的这种转换机制。

